%% Список основных результатов работы.
%% Подключается в двух местах: Dissertation/conclusion.tex (заключение
%% диссертации) и Synopsis/content.tex (заключение автореферата).
%% Согласно ГОСТ Р 7.0.11-2011, п. 5.3.3 и 9.2.3.
\begin{enumerate}
    \item Установлено, что стандартная специальная обертывающая алгебры Ли
          $N\Phi(\mathbb{K})$ существует для всех типов $\Phi$, кроме типов $D_n$ $(n\geq
          4)$ и $E_n$ $(n=6,7,8)$; для типа $F_4$ построены как стандартные, так и
          нестандартные специальные обертывающие алгебры.

    \item Для алгебр Ли $N\Phi(\mathbb{K})$ классических типов найдены условия
          однозначности специальных обертывающих алгебр $R_{\Phi}$. Специальная
          обертывающая алгебра типа $A_{n-1}$ $(n>3)$ стандартна тогда и только тогда,
          когда она или ее противоположная алгебра по модулю аннулятора изоморфна
          $NT(n,\mathbb{K})/\mathbb{K}e_{n1}$; ассоциативная среди стандартных
          единственна и изоморфна $NT(n,\mathbb{K})$. Для типа $B_n$ среди
          стандартных специальных обертывающих алгебр, у которых
          $R_{\Phi}/T_{2,-1}\simeq NT(n+1,\mathbb{K})$, условие монотонности
          "--- вложение каждого члена последовательности подалгеброй в
          следующий "--- однозначно выделяет алгебру $RB_n(\mathbb{K})$.

    \item Доказано, что алгебра $NT(n+1,\mathbb{K})$ "--- единственная с точностью до
          изоморфизма ассоциативная точная обертывающая алгебры Ли $N\Phi(\mathbb{K})$
          типа $A_n$ $(n\geq 3)$, допускающая градуировку, согласованную с градуировкой
          Картана. Условие стандартности такой алгебры слабее и однозначности уже не
          дает.

    \item Для алгебр Ли $N\Phi(\mathbb{K})$ типов $B_n$, $C_n$ и $D_n$ над произвольным
          полем построены левосимметрические точные обертывающие алгебры в виде
          полупрямых сумм алгебры $NT(n,\mathbb{K})$ и подходящего идеала алгебры Ли
          $N\Phi(\mathbb{K})$.

    \item Решена проблема~1 из~\cite{sigsam2001}: найдено число стандартных идеалов
          алгебры Ли $N\Phi(q)$ над конечным полем для всех классических и исключительных
          типов $\Phi$. Перечислены нестандартные идеалы обертывающей алгебры $RD_n(q)$,
          чем завершено перечисление идеалов специальных обертывающих алгебр $R_\Phi$
          классических типов; при этом впервые вычислена трехкратная комбинаторная сумма
          с $q$-биномиальными коэффициентами методом интегрального представления
          комбинаторных сумм.
\end{enumerate}
